% ===============================
% Worksheet / Manual Template
% ===============================
\documentclass[12pt, a4paper]{article}

% ===============================
% Metadata
% ===============================
\newcommand*{\CourseCode}{MAT215}
\newcommand*{\CourseTitle}{Complex Variables \& Laplace Transform}
\newcommand*{\Chapter}{Laplace Transform of Step Functions}
\newcommand*{\Topics}{The Second Translation Formula}
\newcommand*{\DocDate}{April 18, 2025}

\include{header.tex}

% ==========================================
% Document
% ==========================================
\begin{document}
\FrontPage
\Large


\begin{heading}
	What is Step Function?
\end{heading}

\vspace{10pt}

\begin{definition}[The Unit Step Function]
	The unit step function, denoted by $u(t)$, is defined as
	\[
		u(t) = \begin{cases}
			0, & t < 0, \\
			1, & t \geq 0.
		\end{cases}
	\]
\end{definition}

\vspace{10pt}

\begin{definition}[Shifted Unit Step Function]
	The unit step function, denoted by $u(t-a)$, is defined as
	\[
		u(t-a) = \begin{cases}
			0, & t < a, \\
			1, & t \geq a.
		\end{cases}
	\]
\end{definition}

\clearpage

\begin{heading}
	Physical Interpretation of Unit Step Function
\end{heading}

\clearpage

\begin{heading}
	The second translation formula
\end{heading}

\vspace{10pt}

\begin{definition}[Laplace Transform of $f(t) \cdot u(t-a)$]
	If the Laplace transform of $f(t)$ exists, then the Laplace transform of $f(t) \cdot u(t-a)$ is given by
	\[
		\mathcal{L}\{f(t) \cdot u(t-a)\} = e^{-as}\mathcal{L}\{f(t+a)\},
	\]
	where $a > 0$.
\end{definition}

\clearpage

\begin{heading}
	Important Trigonometric Identities
\end{heading}

\vspace{10pt}

\begin{definition}[Compound Angle Identities]
	\begin{align*}
		\sin(A + B) &= \sin A \cos B + \cos A \sin B, \\
		\sin(A - B) &= \sin A \cos B - \cos A \sin B, \\
		\cos(A + B) &= \cos A \cos B - \sin A \sin B.\\
		\cos(A - B) &= \cos A \cos B + \sin A \sin B.
	\end{align*}
\end{definition}

\vspace{10pt}

\begin{definition}[Product-to-Sum Identities]
	\begin{align*}
		2 \sin A \cos B &= \sin(A+B) + \sin(A-B), \\
		2 \cos A \cos B &= \cos(A+B) + \cos(A-B), \\
		2 \sin A \sin B &= \cos(A-B) - \cos(A+B).
	\end{align*}
\end{definition}

\vspace{10pt}

\begin{definition}[Negative Angle Identities]
	\begin{align*}
		\sin(-A) &= -\sin(A), \\
		\cos(-A) &= \cos(A).
	\end{align*}
\end{definition}

\clearpage

\begin{heading}
	How to compute $f(t+a)$?
\end{heading}

\clearpage

\begin{heading}
	Application of the second translation formula
\end{heading}

\vspace{10pt}

\begin{problem}
	Evaluate 
	\[\mathcal{L}\{(2t-3) \ u(t-1)\}\]
\end{problem}

\clearpage

\begin{problem}
	Evaluate 
	\[\mathcal{L}\{e^{-2t} \ u(t-1)\}\]
\end{problem}

\clearpage

\begin{problem}
	Evaluate 
	\[\mathcal{L}\{t^2 e^{2t} \ u(t-2)\}\]
\end{problem}

\clearpage

\begin{problem}
	Evaluate 
	\[\mathcal{L}\{t \sin(2t) \ u(t-\pi)\}\]
\end{problem}

\clearpage

\begin{problem}
	Evaluate 
	\[\mathcal{L}\{e^{-2t} \cos t\ u(t-3\pi)\}\]
\end{problem}

\clearpage

\begin{problem}
	Evaluate 
	\[\mathcal{L}\{t e^{3t} \sin t\ u(t-\pi)\}\]
\end{problem}

\clearpage

\begin{heading}
	Converting a piecewise function into Step function
\end{heading}

\clearpage

\begin{heading}
	Laplace Transform of Piecewise Functions
\end{heading}

\vspace{10pt}

\begin{problem}
	Find the Laplace transform of the function
	\[
	f(t)=
	\begin{cases}
	0, & 0\le t<2,\\
	3, & 2\le t<4,\\
	e^t, & t\ge 4.
	\end{cases}
	\]
\end{problem}

\clearpage

\begin{problem}
	Find the Laplace transform of the function
	\[
	f(t)=
	\begin{cases}
	0, & 0\le t<\pi,\\
	\cos t, & t\ge \pi.
	\end{cases}
	\]
\end{problem}

\clearpage

\begin{problem}
	Find the Laplace transform of the function
	\[
	f(t)=
	\begin{cases}
	5\sin t, & 0\le t<\pi,\\
	0, & t\ge \pi.
	\end{cases}
	\]
\end{problem}

\clearpage

\begin{problem}
	Find the Laplace transform of the function
	\[
	f(t)=
	\begin{cases}
	0, & 0\le t<\pi,\\
	\sin t, & \pi\le t<2\pi,\\
	0, & t\ge 2\pi.
	\end{cases}
	\]
\end{problem}

\clearpage

\begin{problem}
	Find the Laplace transform of the function
	\[
	f(t)=
	\begin{cases}
	5\sin t, & 0\le t<\pi,\\
	-4\cos(2t), & t\ge \pi.
	\end{cases}
	\]
\end{problem}

\clearpage



\EndPage
\end{document}